% ============================================
% Results
% ============================================
\section{Results}
\label{sec:results}

% --------------------------------------------
\subsection{Community-level selection is prevalent in microbial coalescence}
\label{sec:results:cls_prevalent}

Our central question is whether selection acts primarily on whole communities as cohesive units during community coalescence. To address this, we asked how often coalescence outcomes are represented primarily by one parental community versus blending of both or forming a novel community. Consider a coalescence event where parental communities A and B mix to produce community C (\figref{fig:fig1}A). We represent each community by its normalized abundance vector and quantify the cosine similarity between the coalesced community C and each parent (\figref{fig:fig1}B):
\begin{equation}
\text{Sim}(C, A) = \vec{x}_C \cdot \vec{x}_A, \quad \text{Sim}(C, B) = \vec{x}_C \cdot \vec{x}_B
\label{eq:similarity}
\end{equation}
where $\vec{x}_A$, $\vec{x}_B$ are normalized abundance vectors of parents A and B, and $\vec{x}_C$ is that of the coalesced community. These two similarity scores yield a two-dimensional similarity map in which the location of C reflects how strongly it resembles each parent. We partition this map into three classes, each of which corresponds to qualitatively different types of coalescence outcomes (see Methods). If C lies close to one parent and far from the other, we infer community-level selection (CLS) for that parental community. If C lies near the diagonal and similarly resembles both parents, it is consistent with a Mixture where both parental communities coexist. If C lies far from both parents, the community has reassembled into a novel configuration, which we term Restructuring.

To quantify the occurrence of outcome types experimentally, we assembled random in-vitro communities and subjected them to pairwise coalescence (\figref{fig:fig1}C). We first curated a strain library of 54 bacterial isolates collected from diverse environments (soil, tree surface, and flower stamen). The library is phylogenetically broad, spanning 29 families across three phyla: Proteobacteria, Firmicutes, and Bacteroidota (Supplementary Figs. 1, 2). From this library, we assembled 42 parental communities with varying initial richness (6, 12, or 24 species; see Methods) and stabilized them for 7 days under daily growth--dilution cycles ($\times$30 every 24~h) (\figref{fig:fig1}C). Our initial experiments were done in Base medium (1~g~L$^{-1}$ yeast extract, 1~g~L$^{-1}$ soytone, 10~mM sodium phosphate, 0.1~mM CaCl$_2$, 2~mM MgCl$_2$, 4~mg~L$^{-1}$ NiSO$_4$, 50~mg~L$^{-1}$ MnCl$_2$, 5~g~L$^{-1}$ glucose and 4~g~L$^{-1}$ urea; pH~6.5), a buffered complex medium that we have previously determined has moderate interaction strength and supports high species coexistence (Supplementary Information, mean species survival ratio = $74 \pm 2$\%). We then performed 83 pairwise coalescence events by mixing parents at equal volume and restabilizing for an additional 7 days. Community composition of parental communities and coalesced communities were measured at the end of stabilization by 16S rRNA amplicon sequencing (Amplicon Sequence Variants, ASVs), and we also recorded the biomass (OD) and pH of the communities to contextualize assembly and post-coalescence dynamics.

Representative time series illustrate the spectrum of post-coalescence dynamics (\figref{fig:fig1}D). In outcomes classified as CLS, one parental lineage rapidly displaces the other after mixing and the coalesced community converges to a composition closely matching that parent, while largely preserving its internal relative-abundance structure over serial dilution cycles. This behavior is consistent with selection acting at the community level rather than on individual species alone. In the Restructuring case, the merged community converges toward a novel state distinct from both parental communities. Among 12 randomly selected time-series trajectories, we did not observe Mixture cases in which both parents remained comparably represented after coalescence (see Supplementary Fig. 3).

Projecting all outcomes into the similarity space (\figref{fig:fig1}E) revealed that CLS is the most frequent empirical outcome. Of the 82 analyzable coalescence events (one excluded due to low sequencing depth), 54 were classified as CLS, 26 as Restructuring, and only 2 as Mixture (\figref{fig:fig1}E). This pattern of CLS as the most frequent outcome was robust across variants of similarity metrics (Supplementary Fig. 4). One possibility is that the observed frequency of CLS results from skewed species abundance distributions rather than correlated selection among species from the same parental community. To rule out this possibility, we compared experimental outcomes against two null models that assume no correlation in species selection: (1) abundance-weighted random selection and (2) shuffled abundance. The experimentally observed asymmetry in CLS significantly exceeded both null expectations (Supplementary Fig. 5), indicating that CLS reflects genuine ecological dynamics beyond statistical artifacts of skewed abundance distributions.

% Figure 1
\begin{figure}[H]
\centering
\includegraphics[width=\textwidth]{figure_main_1_experimental_overview.pdf}
\caption{\textbf{Community-level selection is the dominant outcome in experimental microbial coalescence.}
\textbf{(A)} Coalescence schematic: communities A and B mix to produce C, classified as CLS, Mixture, or Restructuring.
\textbf{(B)} Cosine similarity quantifies compositional resemblance.
\textbf{(C)} Experimental workflow: 54 isolates assembled into 42 parental communities, pairwise coalesced, tracked by 16S sequencing.
\textbf{(D)} Representative time-series: two CLS outcomes and one Restructuring outcome.
\textbf{(E)} CLS is the dominant outcome (65.9\%, 54/82; $n = 82$ coalescence events), followed by Restructuring (31.7\%) and Mixture (2.4\%).}
\label{fig:fig1}
\end{figure}


% --------------------------------------------
\subsection{Theoretical model with random gLV interactions reproduces community-level selection}
\label{sec:results:glv_model}
To gain insight into why CLS is the prevalent outcome, we built a generalized Lotka--Volterra (gLV) model that mirrors the experimental protocol (\figref{fig:fig2}A). In this classic ecological framework, species grow logistically and compete pairwise:
\begin{equation}
\frac{dN_i}{dt} = g_i N_i \left(1 - \frac{1}{k_i} \sum_j \alpha_{ij} N_j \right)
\label{eq:glv}
\end{equation}
where $N_i$ is abundance, $g_i$ is growth rate, $k_i$ is carrying capacity, and $\alpha_{ij}$ is the competition coefficient between species $i$ and $j$. We fixed $g_i = k_i = 1$ and drew competition coefficients randomly with mean $\mu$, which controls the average competition strength: higher $\mu$ means stronger interspecies competition \citep{Hu2022,Hu2025}. From a pool of 48 species, we generated two parental communities of 12 species each with no shared species, allowed each to reach equilibrium, then mixed them equally to simulate coalescence (\figref{fig:fig2}A; see Methods for details). Post-coalescence compositions were analyzed using the same similarity metrics as in the experiments.

We simulated 200 random coalescence events at interaction strength $\mu = 0.6$ and analyzed the similarity of each coalescence outcome to the two parental communities. At this interaction strength, the random-interaction model quantitatively reproduces the high frequency of CLS observed in the experiments. The distribution of outcomes in the similarity map concentrates near the A-like or B-like axes, indicating that one parental community is dominantly selected rather than forming symmetric mixtures (\figref{fig:fig2}B). Quantitatively, CLS accounts for 61\% of all outcomes, far more frequent than Restructuring (26\%) or Mixture (13\%; \figref{fig:fig2}B, right). This observation suggests that interspecies interactions alone, a minimal and generic condition, are sufficient to reproduce community-level selection in coalescence. 

To understand how CLS emerges in the model, we examined the role of the assembly process. During assembly, competitive exclusion filters out species with strong mutual interactions, leaving communities with reduced mean interaction strength compared to the initial pool (paired $t$-test, $p < 0.001$; \figref{fig:fig2}C; Supplementary Fig. 26)\cite{Gilpin1994, Lechon2021}. This generates internal coherence, whereby species that survived assembly together exhibit weak competitive interactions and consequently share survival or extinction during coalescence. \jy{Theoretical work on colonization resistance (the ability to resist invasion) has shown that communities with weak or facilitative interactions among residents can exhibit stronger resistance to invasion\citep{Tikhonov2016, Niehaus2021}, and recent experiments confirm that interspecies interaction strength is a key determinant of community susceptibility to invaders\citep{OrozcoFuentes2024}.} Consistent with this mechanism, species pairs from the same parent showed strongly positive selection correlation across coalescence events, whereas cross-parent pairs showed negative correlation (\figref{fig:fig2}D). This pattern appears in both simulation and experiment, confirming that assembly-driven coherence underlies CLS even when all pairwise interactions are drawn from a simple random distribution.

% Figure 2
\begin{figure}[H]
\centering
\includegraphics[width=\textwidth]{figure_main_2_simulation_framework.pdf}
\caption{\textbf{Random competitive interactions are sufficient to reproduce community-level selection.}
\textbf{(A)} gLV simulation workflow: competition coefficients drawn randomly with mean competition strength $\mu$; communities (48 total species in pool, 12 species per community) assembled to steady state, mixed, restabilized.
\textbf{(B)} At $\mu = 0.6$, simulations reproduce high CLS ($n = 200$).
\textbf{(C)} Assembly reduces mean interaction strength ($n = 200$; paired $t$-test, $p < 0.001$).
\textbf{(D)} Same-parent species show positive selection correlation; cross-parent pairs show negative correlation (simulation: $\mu = 0.6$, $n = 200$; experiment: $n = 82$; permutation test, $p < 0.001$).}
\label{fig:fig2}
\end{figure}


% --------------------------------------------
\subsection{Interaction strength controls the occurrence of community-level selection}
\label{sec:results:interaction_strength}

Having established that random interspecies interactions can recapitulate community-level selection, we next investigated how interaction strength alters coalescence outcomes. We simulated 1,000 coalescence events at each of three representative interaction strengths ($\mu = 0.3, 0.6, 0.8$) and mapped outcomes into the similarity space (\figref{fig:fig3}A). At weak interactions ($\mu = 0.3$), outcomes clustered in the interior of the map, indicating that both parental communities persisted comparably after mixing. As $\mu$ increased to 0.6 and 0.8, outcomes shifted toward the axes, indicating that one parent increasingly dominated the other.

To quantify this shift, we defined the Parental Dominance Index (PDI), which measures how strongly the coalesced community favors one parent over the other. PDI ranges from 0 (complete dominance by parent B) through 0.5 (equal contributions from both parents) to 1 (complete dominance by parent A). At $\mu = 0.3$, PDI values concentrated near 0.5, consistent with Mixture-dominated outcomes. At $\mu = 0.6$ and $0.8$, PDI distributions became bimodal and skewed toward 0 or 1, reflecting more frequent and stronger CLS (\figref{fig:fig3}A, histograms).

Across interaction strengths $\mu = 0$–$1.2$, we observed a qualitative transition in coalescence outcomes (\figref{fig:fig3}B). At low $\mu$, Mixture was the most frequent outcome and CLS was rare. As $\mu$ increased, the Mixture fraction declined while CLS rose and eventually dominated. Restructuring emerged at modest to high interaction strengths. These patterns were robust to variation in carrying capacities, interaction distributions, similarity metrics, and community size (Supplementary Figs. 6–9). \jy{Notably, the frequency of CLS remained relatively stable across parental communities ranging from 4 to 48 species, spanning our experimental species pool range.} The pairwise selection correlation also increased with interaction strength, indicating that amplification of the assembly effect underlies the emergence of CLS (Supplementary Fig. 10). Together, these simulation results imply a bidirectional shift: weakening interactions drives outcomes toward Mixture, whereas strengthening interactions drives them toward CLS.

% Figure 3
\begin{figure}[H]
\centering
\includegraphics[width=\textwidth]{figure_main_3_interaction_strength.pdf}
\caption{\textbf{Stronger interactions shift coalescence outcomes from Mixture toward community-level selection.}
\textbf{(A)} Coalescence-outcome maps at three interaction strengths ($\mu = 0.3, 0.6, 0.8$; $n = 1{,}000$ coalescence events per condition). Density maps show outcome distribution in similarity space; histograms show Parental Dominance Index (PDI) distributions. Weak interactions ($\mu = 0.3$) produce Mixture-dominated outcomes; stronger interactions shift outcomes toward CLS.
\textbf{(B)} Outcome fractions across interaction strengths ($\mu = 0$--$1.2$; $n = 1{,}000$ per $\mu$ value). Increasing $\mu$ shifts outcomes from Mixture toward CLS.}
\label{fig:fig3}
\end{figure}


% --------------------------------------------
\subsection{Nutrient-dependent interaction strength in experiments recapitulates model predictions}
\label{sec:results:nutrient_experiment}

We asked whether the model-predicted shift in coalescence outcomes depending on interaction strength can be recapitulated experimentally. Following prior work showing that nutrient availability intensifies microbial competition \citep{Hu2022, Ratzke2020, Hu2025}, we conducted additional coalescence experiments by removing or augmenting glucose and urea in the Base medium used in \figref{fig:fig1} (Methods). This yielded two additional media conditions (\figref{fig:fig4}A): Nutr$-$ (no added glucose/urea) and Nutr+ (high supplementation). Higher nutrient availability is expected to amplify consumer--resource feedbacks and intensify environmental modification, thereby strengthening interspecies interactions \citep{Ratzke2020}. To empirically validate this effect, we measured failed invasion frequency using pairwise invasion assays among the 12 most abundant isolates (95:5 initial frequency; Methods, Supplementary Figs. 11–13). The fraction of failed invasions---which indicates strong competitive exclusion---increased monotonically with nutrient supply (\figref{fig:fig4}B). Calibrating these values against gLV simulations yielded approximate mean interaction strengths of $\mu \approx 0.5$ for Nutr$-$, $\mu \approx 0.7$ for Base, and $\mu \approx 0.9$ for Nutr+ (see Supplementary Methods).

We tested this prediction experimentally across all three media using the same parental community library. The distribution of outcomes in similarity space shifted systematically with nutrient concentration (\figref{fig:fig4}C), which we quantified by the fraction of each outcome type (\figref{fig:fig4}D). In Nutr$-$ ($n = 90$), where interactions are weakest, Mixture was the most frequent outcome (56\%), and CLS was substantially reduced to 34\% compared to the Base medium (CLS: 60\%, Mixture: 8\%), indicating that community-level selection largely disappears under weak interactions. In Nutr+ ($n = 90$), CLS further increased to 72\%. Pairwise selection correlation analysis further supported this pattern: in Nutr$-$, species from the same parent showed no significant difference in selection correlation compared to cross-parent pairs, whereas in Nutr+, same-parent species exhibited significantly higher selection correlation than cross-parent pairs ($p < 0.001$; Supplementary Fig. 14), indicating that correlated species selection underlies the increased CLS at higher interaction strengths. These results are consistent with the model prediction that the frequency of CLS increases with the strength of interaction. When interactions are weak, CLS largely disappears and both parental communities mix; in contrast, stronger interactions are dominated by CLS.

% Figure 4
\begin{figure}[H]
\centering
\includegraphics[width=\textwidth]{figure_main_4_nutrient_perturbation.pdf}
\caption{\textbf{Nutrient modulation validates the model-predicted dependence of coalescence outcomes on interaction strength.}
\textbf{(A)} Nutrient supply modulates interactions: Nutr$-$ (no glucose/urea), Base, Nutr+ (high supplementation).
\textbf{(B)} Failed invasion frequency increases with nutrients ($n = 132$ assays per medium), confirming stronger competition.
\textbf{(C)} Coalescence maps and PDI histograms per medium (Nutr$-$: $n = 90$; Base: $n = 82$; Nutr+: $n = 90$). Shapes indicate initial richness.
\textbf{(D)} CLS increases with nutrients (34\% $\rightarrow$ 60\% $\rightarrow$ 72\%); Mixture declines (56\% $\rightarrow$ 8\% $\rightarrow$ 6\%), confirming model predictions.}
\label{fig:fig4}
\end{figure}


% --------------------------------------------
\subsection{Predictability of community-level selection separates into two regimes}
\label{sec:results:predictability}

We observed CLS in both Base and Nutr+ media, and we next asked whether we can predict the direction of CLS—that is, which parental community will be selected during coalescence. Prior work has suggested that species-level competitive outcomes may correlate with community-level selection direction \citep{Rillig2015, Castledine2020}. Inspired by these observations, we tested whether pairwise competition between dominant species \citep{Hu2022, Ratzke2020} predicts the direction of CLS (\figref{fig:fig5}A).

The relative abundance of dominant species increased with nutrient concentration (43.8\% in Nutr$-$, 51.1\% in Base, 66.9\% in Nutr+; \figref{fig:fig5}B), suggesting that dominant species have a larger relative population size under stronger interactions. We therefore tested whether pairwise competition between dominant species predicts the direction of CLS by analyzing the linear relationship between dominant-species competitive success (from invasion assays) and the PDI of coalescence outcomes (\figref{fig:fig5}C). Predictive power varied markedly across media: in Nutr$-$, where CLS is rare, pairwise competition showed no predictive power ($R^2 = 0.00$); in Base medium, predictive power was weak ($R^2 = 0.11$), suggesting that collective multi-species dynamics shape outcomes; in Nutr+, pairwise competition was substantially more predictive ($R^2 = 0.49$), indicating that the dominant species largely determines which community wins. These results reveal two distinct regimes underlying CLS: an emergent regime at intermediate interactions where multi-species dynamics collectively shape the outcome, and a species-driven regime at strong interactions where a few dominant taxa determine the winner. \jy{We explored the mechanistic basis of the species-driven regime by examining environmental modification through pH. In our system, dominant species are often strong pH modifiers that either acidify or alkalinize the medium (Supplementary Fig. 15). Accordingly, the class of dominant species (acidifier vs.\ alkalizer) is highly correlated with parental community pH at stabilization (Supplementary Fig. 16). The predictive power of pH depends on nutrient conditions: in Base medium, acidic communities (pH $<$ 6.5) win only 56\% of matchups against alkaline communities (pH $>$ 7.5; $n = 41$), whereas in Nutr+ medium, acidic communities win 91\% of such matchups ($n = 32$, $p < 0.0001$; Supplementary Fig. 17). This nutrient-dependent pattern suggests that amplified metabolic activity under high nutrients intensifies pH modification and excludes pH-sensitive competitors. Thus, environmental modification via pH contributes to the species-driven regime, where dominant taxa determine coalescence outcomes through strong competitive interactions.}

% Figure 5
\begin{figure}[H]
\centering
\includegraphics[width=\textwidth]{figure_main_5_predictability.pdf}
\caption{\textbf{Two mechanistic regimes underlie community-level selection: emergent versus species-driven dynamics.}
\textbf{(A)} CLS may arise from emergent dynamics (multi-species) or species-driven dynamics (dominant species).
\textbf{(B)} Dominant species abundance increases with nutrients (Nutr$-$: $n = 90$; Base: $n = 82$; Nutr+: $n = 90$). Error bars: s.e.m.
\textbf{(C)} Predictability of CLS direction from dominant-species pairwise competition. Invasion success (x-axis) vs PDI (y-axis). Predictive power increases: Nutr$-$ $R^2 = 0.00$; Base $R^2 = 0.11$ (emergent regime); Nutr+ $R^2 = 0.49$ (species-driven regime). Linear regression; shaded regions: 95\% CI.}
\label{fig:fig5}
\end{figure}


% --------------------------------------------
\jy{\subsection{Community-level selection generalizes to natural sample-derived communities}
\label{sec:results:natural_communities}}

\jy{The synthetic communities described above were assembled from individually isolated bacterial strains, allowing precise control over initial composition but raising the question of whether CLS generalizes to more ecologically realistic settings. To address this, we performed coalescence experiments using communities derived from natural environmental samples, which harbor complex species assemblages shaped by evolutionary and ecological processes in their native habitats.}

\jy{We collected six environmental samples from diverse microhabitats (soil, compost, and decomposing organic matter) and established bacterial communities through seven rounds of serial growth-dilution in laboratory media, allowing natural assemblages to stabilize under controlled conditions (\figref{fig:fig6}A). After stabilization, we performed 16 pairwise coalescence events across all three nutrient conditions (Nutr$-$, Base, and Nutr+) and analyzed outcomes using the same similarity framework applied to synthetic communities. The results revealed striking parallels between natural and synthetic systems (\figref{fig:fig6}B,C). Natural sample-derived communities exhibited the same qualitative patterns: outcomes clustered along the axes in similarity space, indicating one-sided selection dynamics rather than symmetric mixing. Moreover, the fraction of CLS outcomes increased with nutrient availability, mirroring the interaction-strength dependence observed in synthetic communities. These results demonstrate that community-level selection is a robust phenomenon that emerges regardless of community assembly history---from precisely controlled synthetic consortia to complex, naturally-evolved microbial assemblages---suggesting that the underlying mechanism of correlated species retention operates broadly across microbial systems.}

% Figure 6
\begin{figure}[H]
\centering
\includegraphics[width=\textwidth]{figures/figure_main_6_natural_communities.pdf}
\caption{\textbf{Coalescence outcomes in natural sample-derived communities.}
\textbf{(A)} Experimental workflow: natural environmental samples are cultured through serial growth-dilution to establish stable communities, then mixed pairwise for coalescence experiments.
\textbf{(B)} Phase diagrams showing coalescence outcomes in parental similarity space across three nutrient conditions. Each point represents a coalescence event; dashed lines indicate equal similarity to both parental communities. Histograms below show the distribution of Parental Dominance Index (PDI).
\textbf{(C)} Stacked bar chart summarizing coalescence outcome fractions (CLS: community-level selection, M: Mixture, R: Restructuring) across media conditions. Natural communities exhibit the same CLS-dominated patterns observed in synthetic communities, demonstrating the generality of one-sided selection dynamics.}
\label{fig:fig6}
\end{figure}
